\section{Parallel composition}

The definitions below apply to any two LTSs whose transitions
$s \step{\lambda}{c} s'$ carry a choice $c \in \{\epsilon, t, (t,t')\}$, where
$\epsilon$ is the empty choice and is not written: the memory LTS only has
transitions with choice $\epsilon$. Choices compose by
$\epsilon \cdot c = c \cdot \epsilon = c$ and $t \cdot t' = \{t,t'\}$.

\begin{definition}[Synchronisation]\label{def:sync}
A \emph{synchronisation predicate} between two LTSs with external actions
$A_1$ and $A_2$ is a relation ${\sync} \subseteq A_1 \times A_2$. Between the
memory and the programs, $\syncM$ is the least relation such that
\[
\memop{\sop{E}_U} \syncM \req{U}
\qquad\quad
\memop{\sop{E}_{\mathsf{init}}} \syncM \initp
\qquad\quad
\msend{m} \syncM \mrecv{m}.
\]
Between two programs, $\syncC$ relates the dual actions on a channel:
$c!v \syncC c?v$ and $c?v \syncC c!v$.
\end{definition}

\begin{definition}[Parallel composition of LTSs]\label{def:par}
Given a synchronisation predicate $\mathsf{sync}$, the parallel composition
$L_1 \lpar_{\mathsf{sync}} L_2$ of two LTSs $L_1$ and $L_2$ is the LTS whose
states are the pairs $s_1 \lpar s_2$ of a state of $L_1$ and a state of $L_2$,
whose external actions are $A_1 \uplus A_2$, and whose transitions are given
by the rules
\begin{mathpar}
\inferrule*[left=Left]
  {s_1 \step{\lambda}{c} s_1'}
  {s_1 \lpar s_2 \step{\lambda}{c} s_1' \lpar s_2}

\inferrule*[left=Right]
  {s_2 \step{\lambda}{c} s_2'}
  {s_1 \lpar s_2 \step{\lambda}{c} s_1 \lpar s_2'}

\inferrule*[left=Sync]
  {s_1 \step{\mu_1}{c_1} s_1' \\ s_2 \step{\mu_2}{c_2} s_2' \\ \mu_1 \sync \mu_2}
  {s_1 \lpar s_2 \step{\tau}{c_1 \cdot c_2} s_1' \lpar s_2'}
\end{mathpar}
The symbol $\lpar$ is the composition of LTSs; $\parallel$ remains the
parallel composition of the syntax.
\end{definition}

For a memory $\rho$ and a program $P$, rule Sync gives
$\rho \lpar P \step{\tau}{t} \rho' \lpar P'$ whenever
$\rho \mto{\mu_1} \rho'$, $P \step{\mu_2}{t} P'$ and $\mu_1 \syncM \mu_2$.

\begin{lemma}[Pre-harmony]\label{lem:pre-harmony}
For all $\rho, P, s, \rho', P'$,
\[
\conf{\rho}{P} \step{\tau}{s} \conf{\rho'}{P'}
\iff
\rho \lpar P \step{\tau}{s} \rho' \lpar P'
\]
where $\lpar$ is the composition for $\syncM$ of the memory LTS
(\cref{fig:mem}) and the program LTS (\cref{fig:prog}).
\end{lemma}

That is, the LTS of lqCCS is decomposed into the LTS of programs, which runs
without the quantum memory (\cref{fig:prog}), and the LTS of the memory
(\cref{fig:mem}).

\begin{proof}
($\Leftarrow$) We go by case analysis on the derivation of
$\rho \lpar P \step{\tau}{s} \rho' \lpar P'$.
The only two cases are an internal transition of the program (rule Right) or
a synchronisation (rule Sync), since the memory has no internal transition.
In both cases, we go by induction on the transition of the program
$P \step{\lambda}{s'} P'$: a channel action gives the same transition of the
configuration, with the same memory; a request $\req{U}$, $\initp$ or $\mrecv{m}$
synchronised with the memory gives the rules Op, Init and Meas of
\cref{fig:conf}, since $\sapp{\sop{E}_U} = \upd{U}$ and
$\sapp{\sop{E}_{\mathsf{init}}} = \inirho$.

($\Rightarrow$) For this side, we prove a stronger lemma by generalising the
internal label $\tau$ with any label $\lambda$:
$$
  \conf{\rho}{P} \step{\lambda}{s} \conf{\rho'}{P'}
\Rightarrow
\rho \lpar P \step{\lambda}{s} \rho' \lpar P'
$$
We go by induction on the derivation of
$\conf{\rho}{P} \step{\lambda}{s} \conf{\rho'}{P'}$ and we conclude by induction
hypothesis. The rules Op, Meas and Init of \cref{fig:conf} are instances of the rule
Sync, with the choice $\epsilon \cdot t = t$; the rule Rec is an instance of
the rule Right.
\end{proof}

The configuration of a parallel composition $P \parallel Q$ is itself a
composition of LTSs: the system $\rho \lpar P$ is composed with the program
$Q$, which synchronises with the memory of the system (through $\syncM$) and
with the channels of $P$ (through $\syncC$). Let
${\syncS} \subseteq (A_{\M} \uplus A_{\mathsf{prog}}) \times A_{\mathsf{prog}}$
be the least relation such that
\[
\mu_{\M} \syncS \mu \ \text{ if } \mu_{\M} \syncM \mu
\qquad\qquad
\mu_1 \syncS \mu_2 \ \text{ if } \mu_1 \syncC \mu_2 .
\]

\begin{lemma}[Parallel harmony]\label{lem:par-harmony}
For all $\rho, P, Q, s, \rho', R$,
\[
\conf{\rho}{P \parallel Q} \step{\tau}{s} \conf{\rho'}{R}
\iff
R = P' \parallel Q' \text{ and }
(\rho \lpar P) \lpar Q \step{\tau}{s} (\rho' \lpar P') \lpar Q'
\text{ for some } P', Q',
\]
where the outer composition is for $\syncS$ and the inner one for $\syncM$.
\end{lemma}

\begin{proof}
($\Rightarrow$) By case analysis on the last rule of the transition of
$\conf{\rho}{P \parallel Q}$. ParL is the rule Left, by
\cref{lem:pre-harmony}. ParR is either the rule Right, for a silent move of
$Q$, or the rule Sync between a move of the memory (rule Left inside
$\rho \lpar P$) and a request of $Q$, with choice $\epsilon \cdot t = t$.
SynchL and SynchR are the rule Sync between a channel action of $P$ (rule
Right inside $\rho \lpar P$, which does not change $\rho$) and the dual action
of $Q$, with choice $t \cdot t' = \{t,t'\}$.

($\Leftarrow$) By case analysis on the last rule of the transition of
$(\rho \lpar P) \lpar Q$, the same correspondence in the other direction;
the channel actions of programs always carry a single tag, so that the choice
of a synchronisation between $P$ and $Q$ is a pair $\{t,t'\}$.
\end{proof}
